\section{核心内容}

\begin{overviewbox}
\begin{itemize}
  \item 理解 GPU 架构与组织的基础；
  \item 理解线程调度、延迟容忍与硬件占用率；
  \item 区分驻留线程束与当前实际发射指令的线程束。
\end{itemize}
\end{overviewbox}

\subsection{软件执行层次}

CUDA 的软件执行层次由网格、线程块和线程组成。网格中的所有线程运行同一份核函数代码，每个线程通过线程索引和线程块索引选择自身负责的数据与工作。这种执行方式属于\term{单程序多数据}{Single Program, Multiple Data, SPMD}模型。

线程块可以组织成一维、二维或三维形状。多维组织使线程索引能够自然对应多维数据，从而简化地址计算；维度本身并不会自动提高访存效率或计算性能。例如，二维图像既可以使用二维线程块，也可以使用一维线程块，后一种方式需要显式完成一维索引到二维坐标的映射。

\begin{keybox}
线程块的形状用于表达问题结构；硬件执行方式取决于线程块如何驻留到流式多处理器、如何拆分为线程束，以及线程束如何被调度。
\end{keybox}

\subsection{线程块的三个基本语义}

\begin{enumerate}
  \item \textbf{线程块是资源分配与驻留的基本单位。} 同一线程块的全部线程被分配到同一个流式多处理器；一个线程块不会跨越多个流式多处理器执行。
  \item \textbf{同一线程块内部可以高效协作。} 线程可以使用共享内存，并通过线程块级屏障进行同步。
  \item \textbf{不同线程块的执行顺序没有保证。} 核函数必须在不依赖线程块调度顺序的前提下得到正确结果；同一程序由此能够在硬件并行度不同的 GPU 上透明伸缩。
\end{enumerate}

\subsection{软件层次到硬件层次的映射}

软件执行层次与硬件执行资源之间的基本关系为
\[
\boxed{
\text{网格} \rightarrow \text{线程块} \rightarrow \text{流式多处理器驻留}
\rightarrow \text{线程束} \rightarrow \text{处理块发射与执行}
}.
\]

\term{线程束}{warp}是硬件调度的基本单位。线程块驻留到某个流式多处理器后，其线程按照线性化线程索引划分为若干线程束；截至 Hopper 的所列 NVIDIA 架构均采用 32 线程的线程束。控制发散、延迟隐藏与占用率都以线程束为核心进行分析。
